% !TeX root = part01.tex

\documentclass[../final.tex]{subfiles}

\begin{document}

% ============================================================
% Семинар 1
% ============================================================

\practice{1}{04.09.2026}{Квантовые распределения и вырожденный ферми-газ}


% ============================================================
% Статистические распределения
% ============================================================

\section*{Статистические распределения}

Будем использовать принятую в статистической физике систему единиц,
в которой постоянная Больцмана положена равной единице:
%
\begin{equation*}
    k_B=1.
\end{equation*}
%
Поэтому температура $T$ имеет размерность энергии.
Постоянную Планка при этом сохраняем явно.


% ------------------------------------------------------------
% Микроканоническое распределение
% ------------------------------------------------------------

\subsection*{Микроканоническое распределение}

Рассмотрим изолированную систему с фиксированной энергией $E_0$.
Вероятность обнаружить систему в элементе фазового объёма $d\Gamma$
можно записать в виде
%
\begin{equation*}
    dw
    \sim
    \delta(E-E_0)\,d\Gamma.
\end{equation*}

Дельта-функция показывает, что ненулевой статистический вес
имеют только состояния, удовлетворяющие условию
%
\begin{equation*}
    E=E_0.
\end{equation*}


% ------------------------------------------------------------
% Каноническое распределение
% ------------------------------------------------------------

\subsection*{Каноническое распределение}

Если система находится в тепловом контакте с термостатом температуры
$T$, используется каноническое распределение:
%
\begin{equation*}
    w_k
    =
    \frac{1}{Z}
    e^{-E_k/T},
\end{equation*}
%
где $E_k$ --- энергия $k$-го состояния, а
%
\begin{equation*}
    Z
    =
    \sum_k
    e^{-E_k/T}
\end{equation*}
%
--- статистическая сумма.


\begin{rbox}[left=18pt,right=18pt,top=14pt,bottom=10pt]{[]}
    \centering

    \begin{tikzpicture}[
        >=Latex,
        every node/.style={text=rtext}
    ]

        % Термостат
        \draw[
            draw=rc3,
            line width=1.0pt
        ]
        (0,0) circle (2.15);

        % Система
        \draw[
            draw=rtext,
            line width=0.9pt
        ]
        (0,0) circle (0.75);

        % Небольшой акцент
        \fill[rc3] (0,0) circle (2.2pt);

        \node at (0,0.38) {\small система};
        \node at (0,-0.35) {$E_k$};

        \node at (0,1.55) {$T$};
        \node at (0,-1.55) {\small термостат};

        % Обмен энергией
        \draw[
            <->,
            draw=rtext,
            line width=0.8pt
        ]
        (0.78,0.2)
        .. controls (1.15,0.45) and (1.35,0.55) ..
        (1.55,0.70);

        \node[
            right
        ]
        at (1.18,0.25)
        {\small $E$};

    \end{tikzpicture}

    \captionsetup{
        font=footnotesize,
        labelfont=bf,
        skip=5pt
    }

    \captionof{figure}{
        Система в тепловом контакте с термостатом температуры $T$.
    }
    \label{fig:statphys_canonical}
\end{rbox}


% ------------------------------------------------------------
% Большое каноническое распределение
% ------------------------------------------------------------

\subsection*{Большое каноническое распределение}

Если система может обмениваться с окружающей средой
не только энергией, но и частицами, используется большое
каноническое распределение:
%
\begin{equation*}
    w_{N,k}
    =
    \frac{1}{Z}
    \exp
    \left(
        \frac{\mu N-E_k}{T}
    \right).
\end{equation*}

Здесь
%
\begin{equation*}
    \mu
\end{equation*}
%
--- химический потенциал, а $N$ --- число частиц в системе.


% ============================================================
% Больцмановский газ
% ============================================================

\section*{Условие применимости распределения Больцмана}

Для классического, или больцмановского, газа число частиц
должно быть значительно меньше числа доступных одночастичных
квантовых состояний.

Пусть характерный разброс импульсов частиц равен $\Delta p$.
Число квантовых состояний в фазовом объёме порядка
$V(\Delta p)^3$ оценивается как
%
\begin{equation*}
    N_{\text{сост}}
    \sim
    \frac{
        V(\Delta p)^3
    }{
        (2\pi\hbar)^3
    }.
\end{equation*}

Поэтому условие невырожденности газа имеет вид
%
\begin{equation*}
    \frac{
        V(\Delta p)^3
    }{
        (2\pi\hbar)^3
    }
    \gg
    N.
\end{equation*}

Для теплового движения характерный импульс можно оценить как
%
\begin{equation*}
    \Delta p
    \sim
    mv_T.
\end{equation*}

При этом
%
\begin{equation*}
    v_T
    \sim
    \sqrt{\frac{T}{m}},
\end{equation*}
%
поэтому
%
\begin{equation*}
    \Delta p
    \sim
    \sqrt{mT}.
\end{equation*}

Подставляя эту оценку в условие невырожденности, получаем
%
\begin{equation*}
    \frac{
        V(mT)^{3/2}
    }{
        (2\pi\hbar)^3
    }
    \gg
    N.
\end{equation*}

Или
%
\begin{equation*}
    \frac{
        V m^{3/2}T^{3/2}
    }{
        (2\pi\hbar)^3
    }
    \gg
    N.
\end{equation*}

Отсюда
%
\begin{equation*}
    T^{3/2}
    \gg
    \frac{
        (2\pi\hbar)^3N
    }{
        m^{3/2}V
    }.
\end{equation*}

Возводя обе части в степень $2/3$, получаем
%
\begin{equation*}
    T
    \gg
    \frac{
        (2\pi\hbar)^2
    }{
        mV^{2/3}
    }
    N^{2/3}.
\end{equation*}

Если ввести концентрацию частиц
%
\begin{equation*}
    n
    =
    \frac{N}{V},
\end{equation*}
%
то
%
\begin{equation*}
    T
    \gg
    \frac{
        (2\pi\hbar)^2
    }{
        m
    }
    n^{2/3}.
\end{equation*}

Таким образом, по порядку величины условие применимости
больцмановской статистики можно записать как
%
\begin{equation*}
    T
    \gg
    T_F,
\end{equation*}
%
где $T_F$ --- температура вырождения газа.

Для больцмановского газа среднее число частиц в состоянии $k$
имеет экспоненциальную зависимость от энергии:
%
\begin{equation*}
    \overline{n}_k
    =
    Nw_k
    \sim
    e^{-\varepsilon_k/T}.
\end{equation*}


% ============================================================
% Распределения Ферми и Бозе
% ============================================================

\section*{Распределения Ферми и Бозе}

Для идеального квантового газа среднее число частиц
в одночастичном квантовом состоянии $k$ определяется выражением
%
\begin{equation*}
    n_k
    =
    \left(
        e^{(\varepsilon_k-\mu)/T}
        \pm 1
    \right)^{-1}.
\end{equation*}

Для фермионов выбирается знак ``плюс'':
%
\begin{equation*}
    \boxed{
        n_k^{F}
        =
        \frac{
            1
        }{
            e^{(\varepsilon_k-\mu)/T}+1
        }
    }.
\end{equation*}

Для бозонов выбирается знак ``минус'':
%
\begin{equation*}
    \boxed{
        n_k^{B}
        =
        \frac{
            1
        }{
            e^{(\varepsilon_k-\mu)/T}-1
        }
    }.
\end{equation*}

Для фермионов наличие знака ``плюс'' связано с принципом Паули:
одно одночастичное квантовое состояние не может быть занято
более чем одним фермионом.


% ------------------------------------------------------------
% Ферми-газ при T = 0
% ------------------------------------------------------------

\subsection*{Предел $T\to 0$ для ферми-газа}

Поскольку
%
\begin{equation*}
    C_V
    =
    \left(
        \frac{\partial E}{\partial T}
    \right)_V
    >0,
\end{equation*}
%
при уменьшении температуры энергия системы также уменьшается:
%
\begin{equation*}
    T\downarrow
    \quad\Longrightarrow\quad
    E\downarrow.
\end{equation*}

При абсолютном нуле система находится в основном состоянии:
%
\begin{equation*}
    T=0
    \quad\Longrightarrow\quad
    E=E_{\min}.
\end{equation*}

Для невырожденного основного состояния статистический вес равен
%
\begin{equation*}
    \Gamma=1,
\end{equation*}
%
поэтому
%
\begin{equation*}
    S
    =
    \ln\Gamma
    =
    0.
\end{equation*}


Рассмотрим теперь распределение Ферми--Дирака:
%
\begin{equation*}
    n_k
    =
    \frac{
        1
    }{
        e^{(\varepsilon_k-\mu)/T}+1
    }.
\end{equation*}

Если
%
\begin{equation*}
    \varepsilon_k<\mu,
\end{equation*}
%
то при $T\to0$
%
\begin{equation*}
    \frac{
        \varepsilon_k-\mu
    }{
        T
    }
    \to
    -\infty.
\end{equation*}

Следовательно,
%
\begin{equation*}
    e^{(\varepsilon_k-\mu)/T}
    \to
    0
\end{equation*}
%
и
%
\begin{equation*}
    n_k
    \to
    1.
\end{equation*}

Если же
%
\begin{equation*}
    \varepsilon_k>\mu,
\end{equation*}
%
то
%
\begin{equation*}
    \frac{
        \varepsilon_k-\mu
    }{
        T
    }
    \to
    +\infty,
\end{equation*}
%
следовательно,
%
\begin{equation*}
    e^{(\varepsilon_k-\mu)/T}
    \to
    \infty
\end{equation*}
%
и
%
\begin{equation*}
    n_k
    \to
    0.
\end{equation*}

Таким образом, при $T=0$
%
\begin{equation*}
    n_k
    =
    \begin{cases}
        1,
        &
        \varepsilon_k<\mu,
        \\[4pt]
        0,
        &
        \varepsilon_k>\mu.
    \end{cases}
\end{equation*}

Граничная энергия между заполненными и свободными состояниями
называется \rtxt{3}{энергией Ферми}:
%
\begin{equation*}
    \mu(T=0)
    =
    \varepsilon_F.
\end{equation*}

Поэтому
%
\begin{equation*}
    n_k
    =
    \begin{cases}
        1,
        &
        \varepsilon_k<\varepsilon_F,
        \\[4pt]
        0,
        &
        \varepsilon_k>\varepsilon_F.
    \end{cases}
\end{equation*}


\begin{rbox}[left=18pt,right=18pt,top=14pt,bottom=10pt]{[]}
    \centering

    \begin{tikzpicture}[
        >=Latex,
        every node/.style={text=rtext},
        x=1.15cm,
        y=2.0cm
    ]

        % Оси
        \draw[
            ->,
            draw=rtext,
            line width=0.9pt
        ]
        (0,0) -- (5.1,0)
        node[
            right
        ]
        {$\varepsilon$};

        \draw[
            ->,
            draw=rtext,
            line width=0.9pt
        ]
        (0,0) -- (0,1.45)
        node[
            above
        ]
        {$n(\varepsilon)$};

        % Ступенька Ферми
        \draw[
            draw=rc3,
            line width=1.5pt
        ]
        (0,1)
        -- (3.25,1)
        -- (3.25,0)
        -- (4.75,0);

        % Вспомогательная линия epsilon_F
        \draw[
            draw=rtext,
            line width=0.6pt,
            dashed
        ]
        (3.25,0)
        -- (3.25,1.05);

        % Метки
        \node[
            left
        ]
        at (0,1)
        {$1$};

        \node[
            below
        ]
        at (3.25,0)
        {$\varepsilon_F$};

        \node[
            above right
        ]
        at (1.55,1)
        {$T=0$};

    \end{tikzpicture}

    \captionsetup{
        font=footnotesize,
        labelfont=bf,
        skip=5pt
    }

    \captionof{figure}{
        Распределение Ферми--Дирака при $T=0$:
        все состояния до энергии Ферми заполнены,
        а состояния выше $\varepsilon_F$ свободны.
    }
    \label{fig:statphys_fermi_step}
\end{rbox}


% ============================================================
% Задача 1
% ============================================================

\newpage
\section{Задача 1. Вырожденный электронный газ}

В объёме $V$ находятся $N$ электронов.
Найдём среднюю энергию и давление электронного газа в пределе
%
\begin{equation*}
    T\to0.
\end{equation*}

Электрон имеет две возможные проекции спина, поэтому
спиновая кратность равна
%
\begin{equation*}
    g=2.
\end{equation*}

Для свободного нерелятивистского электрона закон дисперсии имеет вид
%
\begin{equation*}
    \varepsilon
    =
    \frac{p^2}{2m}.
\end{equation*}

При $T=0$ распределение по состояниям представляет собой ступеньку:
%
\begin{equation*}
    n(\varepsilon)
    =
    \begin{cases}
        1,
        &
        \varepsilon<\varepsilon_F,
        \\[4pt]
        0,
        &
        \varepsilon>\varepsilon_F.
    \end{cases}
\end{equation*}


% ------------------------------------------------------------
% Число частиц и энергия Ферми
% ------------------------------------------------------------

\subsection*{Число частиц и энергия Ферми}

Число электронов записывается как интеграл по фазовому пространству:
%
\begin{equation*}
    N
    =
    2
    \int
    \frac{
        d^3p\,d^3r
    }{
        (2\pi\hbar)^3
    }
    n(\varepsilon).
\end{equation*}

Поскольку газ однородный,
%
\begin{equation*}
    \int d^3r
    =
    V,
\end{equation*}
%
и поэтому
%
\begin{equation*}
    N
    =
    2V
    \int
    \frac{
        d^3p
    }{
        (2\pi\hbar)^3
    }
    n(\varepsilon).
\end{equation*}

Перейдём от импульса к энергии.
Из
%
\begin{equation*}
    \varepsilon
    =
    \frac{p^2}{2m}
\end{equation*}
%
следует
%
\begin{equation*}
    p
    =
    \sqrt{2m\varepsilon}.
\end{equation*}

В сферических координатах в импульсном пространстве
%
\begin{equation*}
    d^3p
    =
    4\pi p^2\,dp.
\end{equation*}

Дифференцируя
%
\begin{equation*}
    p
    =
    \sqrt{2m\varepsilon},
\end{equation*}
%
получаем
%
\begin{equation*}
    dp
    =
    \frac{m}{p}\,d\varepsilon.
\end{equation*}

Следовательно,
%
\begin{align*}
    p^2\,dp
    &=
    mp\,d\varepsilon
    \\[3pt]
    &=
    m\sqrt{2m\varepsilon}\,d\varepsilon
    \\[3pt]
    &=
    \frac{
        (2m)^{3/2}
    }{
        2
    }
    \sqrt{\varepsilon}\,d\varepsilon.
\end{align*}

Поэтому
%
\begin{align*}
    d^3p
    &=
    4\pi p^2\,dp
    \\[3pt]
    &=
    2\pi(2m)^{3/2}
    \sqrt{\varepsilon}\,d\varepsilon.
\end{align*}

Таким образом,
%
\begin{equation*}
    \frac{
        d^3p
    }{
        (2\pi\hbar)^3
    }
    =
    \frac{
        (2m)^{3/2}
    }{
        (2\pi\hbar)^3
    }
    2\pi
    \sqrt{\varepsilon}\,d\varepsilon.
\end{equation*}

При $T=0$ состояния заполнены только до энергии Ферми,
поэтому
%
\begin{equation*}
    N
    =
    2V
    \int_0^{\varepsilon_F}
    \frac{
        (2m)^{3/2}
    }{
        (2\pi\hbar)^3
    }
    2\pi
    \sqrt{\varepsilon}\,d\varepsilon.
\end{equation*}

Вынося постоянные множители,
%
\begin{equation*}
    N
    =
    4\pi V
    \frac{
        (2m)^{3/2}
    }{
        (2\pi\hbar)^3
    }
    \int_0^{\varepsilon_F}
    \varepsilon^{1/2}\,d\varepsilon.
\end{equation*}

Интеграл равен
%
\begin{equation*}
    \int_0^{\varepsilon_F}
    \varepsilon^{1/2}\,d\varepsilon
    =
    \frac{2}{3}
    \varepsilon_F^{3/2}.
\end{equation*}

Следовательно,
%
\begin{equation*}
    N
    =
    4\pi V
    \frac{
        (2m)^{3/2}
    }{
        (2\pi\hbar)^3
    }
    \frac{2}{3}
    \varepsilon_F^{3/2}.
\end{equation*}

Отсюда
%
\begin{equation*}
    \varepsilon_F^{3/2}
    =
    \frac{
        3N
    }{
        8\pi V
    }
    \frac{
        (2\pi\hbar)^3
    }{
        (2m)^{3/2}
    }.
\end{equation*}

Возводя обе части в степень $2/3$, получаем
%
\begin{equation*}
    \varepsilon_F
    =
    \left(
        \frac{
            3N
        }{
            8\pi V
        }
    \right)^{2/3}
    \frac{
        (2\pi\hbar)^2
    }{
        2m
    }.
\end{equation*}

После упрощения численного коэффициента:
%
\begin{equation*}
    \boxed{
        \varepsilon_F
        =
        \frac{
            \hbar^2
        }{
            2m
        }
        \left(
            3\pi^2
            \frac{N}{V}
        \right)^{2/3}
    }.
\end{equation*}

Вводя концентрацию
%
\begin{equation*}
    n
    =
    \frac{N}{V},
\end{equation*}
%
получаем
%
\begin{equation*}
    \boxed{
        \varepsilon_F
        =
        \frac{
            \hbar^2
        }{
            2m
        }
        (3\pi^2n)^{2/3}
    }.
\end{equation*}


% ------------------------------------------------------------
% Импульс Ферми
% ------------------------------------------------------------

\subsection*{Импульс Ферми}

По определению
%
\begin{equation*}
    \varepsilon_F
    =
    \frac{
        p_F^2
    }{
        2m
    }.
\end{equation*}

Следовательно,
%
\begin{equation*}
    p_F
    =
    \sqrt{
        2m\varepsilon_F
    }.
\end{equation*}

Подставляя найденную энергию Ферми,
%
\begin{align*}
    p_F
    &=
    \sqrt{
        2m
        \frac{\hbar^2}{2m}
        (3\pi^2n)^{2/3}
    }
    \\[3pt]
    &=
    \hbar
    (3\pi^2n)^{1/3}.
\end{align*}

Таким образом,
%
\begin{equation*}
    \boxed{
        p_F
        =
        \hbar
        (3\pi^2n)^{1/3}
    }.
\end{equation*}

Или через $N$ и $V$:
%
\begin{equation*}
    \boxed{
        p_F
        =
        (3\pi^2)^{1/3}
        \hbar
        \left(
            \frac{N}{V}
        \right)^{1/3}
    }.
\end{equation*}

По порядку величины
%
\begin{equation*}
    p_F
    \sim
    \hbar n^{1/3}.
\end{equation*}


% ------------------------------------------------------------
% Средняя энергия
% ------------------------------------------------------------

\subsection*{Средняя энергия электронов}

Полная энергия электронного газа равна
%
\begin{equation*}
    E
    =
    2
    \int
    \frac{
        d^3p\,d^3r
    }{
        (2\pi\hbar)^3
    }
    \varepsilon
    n(\varepsilon).
\end{equation*}

Для однородного газа
%
\begin{equation*}
    E
    =
    2V
    \int
    \frac{
        d^3p
    }{
        (2\pi\hbar)^3
    }
    \varepsilon
    n(\varepsilon).
\end{equation*}

При $T=0$
%
\begin{equation*}
    E
    =
    2V
    \int_0^{\varepsilon_F}
    \frac{
        (2m)^{3/2}
    }{
        (2\pi\hbar)^3
    }
    2\pi
    \varepsilon^{3/2}\,d\varepsilon.
\end{equation*}

Следовательно,
%
\begin{equation*}
    E
    =
    4\pi V
    \frac{
        (2m)^{3/2}
    }{
        (2\pi\hbar)^3
    }
    \int_0^{\varepsilon_F}
    \varepsilon^{3/2}\,d\varepsilon.
\end{equation*}

Интегрируя,
%
\begin{equation*}
    \int_0^{\varepsilon_F}
    \varepsilon^{3/2}\,d\varepsilon
    =
    \frac{2}{5}
    \varepsilon_F^{5/2},
\end{equation*}
%
получаем
%
\begin{equation*}
    E
    =
    4\pi V
    \frac{
        (2m)^{3/2}
    }{
        (2\pi\hbar)^3
    }
    \frac{2}{5}
    \varepsilon_F^{5/2}.
\end{equation*}

Для числа частиц ранее было найдено
%
\begin{equation*}
    N
    =
    4\pi V
    \frac{
        (2m)^{3/2}
    }{
        (2\pi\hbar)^3
    }
    \frac{2}{3}
    \varepsilon_F^{3/2}.
\end{equation*}

Поэтому
%
\begin{align*}
    \frac{E}{N}
    &=
    \frac{
        \dfrac{2}{5}
        \varepsilon_F^{5/2}
    }{
        \dfrac{2}{3}
        \varepsilon_F^{3/2}
    }
    \\[5pt]
    &=
    \frac{3}{5}
    \varepsilon_F.
\end{align*}

Таким образом, средняя энергия одного электрона равна
%
\begin{equation*}
    \boxed{
        \overline{\varepsilon}
        =
        \frac{E}{N}
        =
        \frac{3}{5}
        \varepsilon_F
    }.
\end{equation*}

Полная энергия газа:
%
\begin{equation*}
    \boxed{
        E
        =
        \frac{3}{5}
        N\varepsilon_F
    }.
\end{equation*}

Подставляя выражение для энергии Ферми,
%
\begin{equation*}
    E
    =
    \frac{3}{5}
    N
    \frac{
        \hbar^2
    }{
        2m
    }
    \left(
        3\pi^2
        \frac{N}{V}
    \right)^{2/3}.
\end{equation*}

Раскрывая степени,
%
\begin{equation*}
    \boxed{
        E
        =
        \frac{
            \hbar^2
        }{
            2m
        }
        \frac{
            3^{5/3}\pi^{4/3}
        }{
            5
        }
        \frac{
            N^{5/3}
        }{
            V^{2/3}
        }
    }.
\end{equation*}

Таким образом,
%
\begin{equation*}
    E
    \propto
    V^{-2/3}
\end{equation*}
%
при фиксированном числе частиц $N$.


% ------------------------------------------------------------
% Давление
% ------------------------------------------------------------

\subsection*{Давление вырожденного электронного газа}

Для классического идеального газа
%
\begin{equation*}
    pV
    =
    NT.
\end{equation*}

Если формально положить $T=0$, эта формула дала бы
%
\begin{equation*}
    p=0.
\end{equation*}

Однако для ферми-газа это неверно.
Даже при абсолютном нуле электроны не могут находиться
в одном и том же одночастичном состоянии из-за принципа Паули.
Поэтому они заполняют состояния с ненулевыми импульсами вплоть
до $p_F$, что приводит к \rtxt{3}{давлению вырождения}.

Используем термодинамическое соотношение
%
\begin{equation*}
    dE
    =
    T\,dS
    -
    p\,dV
    +
    \mu\,dN.
\end{equation*}

Отсюда
%
\begin{equation*}
    p
    =
    -
    \left(
        \frac{
            \partial E
        }{
            \partial V
        }
    \right)_{S,N}.
\end{equation*}

При $T=0$
%
\begin{equation*}
    S=0.
\end{equation*}

Поскольку
%
\begin{equation*}
    E
    =
    C
    N^{5/3}
    V^{-2/3},
\end{equation*}
%
где $C$ не зависит от $V$, имеем
%
\begin{align*}
    \left(
        \frac{
            \partial E
        }{
            \partial V
        }
    \right)_{S,N}
    &=
    -
    \frac{2}{3}
    C
    N^{5/3}
    V^{-5/3}
    \\[3pt]
    &=
    -
    \frac{2}{3}
    \frac{E}{V}.
\end{align*}

Поэтому
%
\begin{equation*}
    \boxed{
        p
        =
        \frac{2}{3}
        \frac{E}{V}
    }.
\end{equation*}

Так как
%
\begin{equation*}
    E
    =
    \frac{3}{5}
    N\varepsilon_F,
\end{equation*}
%
то
%
\begin{align*}
    p
    &=
    \frac{2}{3}
    \frac{1}{V}
    \frac{3}{5}
    N\varepsilon_F
    \\[3pt]
    &=
    \frac{2}{5}
    \frac{N}{V}
    \varepsilon_F.
\end{align*}

Следовательно,
%
\begin{equation*}
    \boxed{
        p
        =
        \frac{2}{5}
        n\varepsilon_F
    }.
\end{equation*}

Подставляя энергию Ферми,
%
\begin{align*}
    p
    &=
    \frac{2}{5}
    n
    \frac{
        \hbar^2
    }{
        2m
    }
    (3\pi^2n)^{2/3}
    \\[3pt]
    &=
    \frac{
        \hbar^2
    }{
        5m
    }
    (3\pi^2)^{2/3}
    n^{5/3}.
\end{align*}

Итак,
%
\begin{equation*}
    \boxed{
        p
        =
        \frac{
            \hbar^2
        }{
            5m
        }
        (3\pi^2)^{2/3}
        n^{5/3}
    }.
\end{equation*}


% ------------------------------------------------------------
% Ответ задачи 1
% ------------------------------------------------------------

\subsection*{Ответ}

Для полностью вырожденного электронного газа
%
\begin{equation*}
    \boxed{
        \varepsilon_F
        =
        \frac{
            \hbar^2
        }{
            2m
        }
        (3\pi^2n)^{2/3}
    },
\end{equation*}
%
%
\begin{equation*}
    \boxed{
        \overline{\varepsilon}
        =
        \frac{3}{5}
        \varepsilon_F
    },
\end{equation*}
%
%
\begin{equation*}
    \boxed{
        p
        =
        \frac{2}{5}
        n\varepsilon_F
        =
        \frac{2}{3}
        \frac{E}{V}
    }.
\end{equation*}


% ============================================================
% Задача 2
% ============================================================

\newpage
\section{Задача 2. Электроны проводимости в меди}

Найти химический потенциал, среднюю энергию и давление
электронов проводимости в меди, если
%
\begin{equation*}
    A
    \simeq
    64\,
    \frac{
        \text{г}
    }{
        \text{моль}
    },
\end{equation*}
%
%
\begin{equation*}
    \rho
    \simeq
    9\,
    \frac{
        \text{г}
    }{
        \text{см}^3
    },
\end{equation*}
%
и каждый атом меди поставляет один электрон проводимости.

В рамках рассматриваемого приближения считаем электронный газ
полностью вырожденным:
%
\begin{equation*}
    T\to0.
\end{equation*}

Поэтому
%
\begin{equation*}
    \mu
    =
    \varepsilon_F.
\end{equation*}


% ------------------------------------------------------------
% Концентрация электронов
% ------------------------------------------------------------

\subsection*{Концентрация электронов проводимости}

В одном моле меди находится число Авогадро атомов:
%
\begin{equation*}
    N_A
    =
    6.022\times10^{23}\,
    \text{моль}^{-1}.
\end{equation*}

Число молей меди в единице объёма равно
%
\begin{equation*}
    \frac{\rho}{A}.
\end{equation*}

Поскольку каждый атом поставляет один электрон проводимости,
концентрация электронов равна
%
\begin{equation*}
    n
    =
    N_A
    \frac{\rho}{A}.
\end{equation*}

Подставляя численные значения,
%
\begin{align*}
    n
    &=
    6.022\times10^{23}
    \frac{
        9
    }{
        64
    }
    \text{см}^{-3}
    \\[3pt]
    &=
    6.022\times10^{23}
    \cdot
    0.140625
    \text{см}^{-3}
    \\[3pt]
    &\simeq
    8.47\times10^{22}
    \text{см}^{-3}.
\end{align*}

Так как
%
\begin{equation*}
    1\,\text{см}^3
    =
    10^{-6}\,\text{м}^3,
\end{equation*}
%
получаем
%
\begin{equation*}
    \boxed{
        n
        \simeq
        8.47\times10^{28}
        \text{м}^{-3}
    }.
\end{equation*}


% ------------------------------------------------------------
% Импульс Ферми
% ------------------------------------------------------------

\subsection*{Импульс Ферми}

Используем выражение
%
\begin{equation*}
    p_F
    =
    \hbar
    (3\pi^2n)^{1/3}.
\end{equation*}

Сначала найдём
%
\begin{align*}
    3\pi^2n
    &\simeq
    29.61
    \cdot
    8.47\times10^{28}
    \\[3pt]
    &\simeq
    2.51\times10^{30}
    \text{м}^{-3}.
\end{align*}

Следовательно,
%
\begin{equation*}
    (3\pi^2n)^{1/3}
    \simeq
    1.36\times10^{10}
    \text{м}^{-1}.
\end{equation*}

Таким образом, ферми-волновое число равно
%
\begin{equation*}
    k_F
    =
    (3\pi^2n)^{1/3}
    \simeq
    1.36\times10^{10}
    \text{м}^{-1}.
\end{equation*}

Используя
%
\begin{equation*}
    \hbar
    =
    1.055\times10^{-34}
    \text{Дж}\cdot\text{с},
\end{equation*}
%
получаем
%
\begin{align*}
    p_F
    &=
    \hbar k_F
    \\[3pt]
    &\simeq
    1.055\times10^{-34}
    \cdot
    1.36\times10^{10}
    \\[3pt]
    &\simeq
    1.43\times10^{-24}
    \frac{
        \text{кг}\cdot\text{м}
    }{
        \text{с}
    }.
\end{align*}

Следовательно,
%
\begin{equation*}
    \boxed{
        p_F
        \simeq
        1.43\times10^{-24}
        \frac{
            \text{кг}\cdot\text{м}
        }{
            \text{с}
        }
    }.
\end{equation*}


% ------------------------------------------------------------
% Химический потенциал
% ------------------------------------------------------------

\subsection*{Химический потенциал}

При $T=0$
%
\begin{equation*}
    \mu
    =
    \varepsilon_F.
\end{equation*}

Энергия Ферми равна
%
\begin{equation*}
    \varepsilon_F
    =
    \frac{
        p_F^2
    }{
        2m_e
    }.
\end{equation*}

Используем массу электрона
%
\begin{equation*}
    m_e
    =
    9.109\times10^{-31}
    \text{кг}.
\end{equation*}

Тогда
%
\begin{align*}
    \varepsilon_F
    &=
    \frac{
        (1.43\times10^{-24})^2
    }{
        2\cdot9.109\times10^{-31}
    }
    \\[5pt]
    &\simeq
    \frac{
        2.05\times10^{-48}
    }{
        1.822\times10^{-30}
    }
    \\[5pt]
    &\simeq
    1.13\times10^{-18}
    \text{Дж}.
\end{align*}

Для перевода в электронвольты используем
%
\begin{equation*}
    1\,\text{эВ}
    =
    1.602\times10^{-19}
    \text{Дж}.
\end{equation*}

Поэтому
%
\begin{align*}
    \varepsilon_F
    &=
    \frac{
        1.13\times10^{-18}
    }{
        1.602\times10^{-19}
    }
    \text{эВ}
    \\[3pt]
    &\simeq
    7.03\,\text{эВ}.
\end{align*}

Следовательно,
%
\begin{equation*}
    \boxed{
        \mu(T=0)
        =
        \varepsilon_F
        \simeq
        7.0\,\text{эВ}
    }.
\end{equation*}


% ------------------------------------------------------------
% Средняя энергия
% ------------------------------------------------------------

\subsection*{Средняя энергия электронов}

Для полностью вырожденного газа
%
\begin{equation*}
    \overline{\varepsilon}
    =
    \frac{3}{5}
    \varepsilon_F.
\end{equation*}

Поэтому
%
\begin{align*}
    \overline{\varepsilon}
    &=
    \frac{3}{5}
    \cdot
    7.03\,\text{эВ}
    \\[3pt]
    &\simeq
    4.22\,\text{эВ}.
\end{align*}

Следовательно,
%
\begin{equation*}
    \boxed{
        \overline{\varepsilon}
        \simeq
        4.22\,\text{эВ}
    }.
\end{equation*}

В джоулях
%
\begin{align*}
    \overline{\varepsilon}
    &=
    \frac{3}{5}
    \cdot
    1.13\times10^{-18}
    \\[3pt]
    &\simeq
    6.76\times10^{-19}
    \text{Дж}.
\end{align*}

То есть
%
\begin{equation*}
    \boxed{
        \overline{\varepsilon}
        \simeq
        6.76\times10^{-19}
        \text{Дж}
    }.
\end{equation*}

Средняя энергия меньше энергии Ферми, поскольку $\varepsilon_F$
является энергией верхнего заполненного состояния, тогда как электроны
занимают все состояния от $0$ до $\varepsilon_F$.


% ------------------------------------------------------------
% Давление
% ------------------------------------------------------------

\subsection*{Давление электронного газа}

Используем найденное в первой задаче соотношение
%
\begin{equation*}
    p
    =
    \frac{2}{5}
    n\varepsilon_F.
\end{equation*}

Подставляя
%
\begin{equation*}
    n
    =
    8.47\times10^{28}
    \text{м}^{-3}
\end{equation*}
%
и
%
\begin{equation*}
    \varepsilon_F
    =
    1.13\times10^{-18}
    \text{Дж},
\end{equation*}
%
получаем
%
\begin{align*}
    p
    &=
    \frac{2}{5}
    \cdot
    8.47\times10^{28}
    \cdot
    1.13\times10^{-18}
    \\[3pt]
    &\simeq
    \frac{2}{5}
    \cdot
    9.54\times10^{10}
    \\[3pt]
    &\simeq
    3.82\times10^{10}
    \text{Па}.
\end{align*}

Следовательно,
%
\begin{equation*}
    \boxed{
        p
        \simeq
        3.82\times10^{10}
        \text{Па}
    }.
\end{equation*}

Это соответствует давлению порядка
%
\begin{equation*}
    p
    \sim
    3.8\times10^5
    \text{атм}.
\end{equation*}

Таким образом, даже при $T\to0$ давление электронного газа
не обращается в нуль.


% ------------------------------------------------------------
% Температура Ферми
% ------------------------------------------------------------

\subsection*{Температура Ферми}

Для оценки степени вырождения электронов меди найдём температуру Ферми:
%
\begin{equation*}
    T_F
    =
    \frac{
        \varepsilon_F
    }{
        k_B
    }.
\end{equation*}

Используя
%
\begin{equation*}
    k_B
    =
    1.381\times10^{-23}
    \frac{
        \text{Дж}
    }{
        \text{К}
    },
\end{equation*}
%
получаем
%
\begin{align*}
    T_F
    &=
    \frac{
        1.13\times10^{-18}
    }{
        1.381\times10^{-23}
    }
    \\[3pt]
    &\simeq
    8.2\times10^4
    \text{К}.
\end{align*}

Следовательно,
%
\begin{equation*}
    \boxed{
        T_F
        \simeq
        8.2\times10^4
        \text{К}
    }.
\end{equation*}

Даже при комнатной температуре $T\simeq300\,\text{К}$
%
\begin{equation*}
    \frac{T}{T_F}
    \simeq
    \frac{
        300
    }{
        8.2\times10^4
    }
    \sim
    4\times10^{-3}
    \ll
    1.
\end{equation*}

Поэтому электроны проводимости меди при обычных температурах
представляют собой сильно вырожденный ферми-газ.


% ------------------------------------------------------------
% Ответ задачи 2
% ------------------------------------------------------------

\subsection*{Ответ}

Для меди в модели свободных электронов получаем
%
\begin{equation*}
    \boxed{
        n
        \simeq
        8.47\times10^{28}
        \text{м}^{-3}
    },
\end{equation*}
%
%
\begin{equation*}
    \boxed{
        \mu
        \simeq
        \varepsilon_F
        \simeq
        7.03\,\text{эВ}
    },
\end{equation*}
%
%
\begin{equation*}
    \boxed{
        \overline{\varepsilon}
        \simeq
        4.22\,\text{эВ}
    },
\end{equation*}
%
%
\begin{equation*}
    \boxed{
        p
        \simeq
        3.82\times10^{10}
        \text{Па}
    }.
\end{equation*}


\end{document}